%% SYNTHETIC TEST UNIT generated by writing/tools/make_outline_book.py from FINAL_OUTLINE_v2 -- not book text
\chapter{光锥排序微扰论与 Drell–Yan 的 Glauber 相消}\label{app:appO}
本单元为测试用的合成文本（由 FINAL\_OUTLINE\_v2 自动生成），陈述内容从略。

\section{第 1 节（测试）}\label{sec:appO-overview}

\begin{definition}[领头区域的数据]\label{def:appO-leading-data}
（测试用合成陈述：内容从略。）
\uses{def:ir-regions, def:method-regions}
\end{definition}

\begin{definition}[光锥时间、中间态与能量分母]\label{def:appO-lcopt}
（测试用合成陈述：内容从略。）
\uses{def:prelim-lightfront}
\end{definition}

\begin{lemma}[由 Feynman 积分导出光锥序微扰论（有限 $\eta_0$） \tagR]\label{lem:appO-lcopt}
（$R_{\mathrm{PT}}$：在重整化微扰论中对红外正规化的部分子态严格。）
（测试用合成陈述：内容从略。）
（本命题只涉及部分子态。）
\uses{def:appO-lcopt, thm:appM-weinberg}
\end{lemma}

\begin{lemma}[分子、瞬时项与瞬时圈 \tagR]\label{lem:appO-numerators}
（$R_{\mathrm{PT}}$：在重整化微扰论中对红外正规化的部分子态严格。）
（测试用合成陈述：内容从略。）
（本命题只涉及部分子态。）
\uses{lem:appO-lcopt}
\end{lemma}

\begin{remark}[与 $V_H$ 重合的接触顶点在领头幂次不出现]\label{rem:appO-contact}
（测试用合成陈述：内容从略。）
% uses: none
\end{remark}

\begin{assumption}[光锥序微扰论的适用性 \tagA]\label{asm:appO-lcopt}
（测试用合成陈述：内容从略。）
地位：测试状态句。
文献：测试。
被使用于：\cref{prop:method-fsi}、\cref{lem:dy-fsi}、\cref{thm:dy-factorization}、\cref{prop:break-measurement}、\cref{lem:appO-hard-locality}、\cref{lem:appO-unitarity}、\cref{prop:appO-ss-localize}、\cref{lem:appO-side-independent-windows}、\cref{prop:appO-ss-cancel}、\cref{prop:appO-twoexch-spectator}、\cref{prop:appO-one-sided}、\cref{thm:appO-main}、\cref{cor:appO-interface}、\cref{prop:appO-tmd}。
解除途径：测试解除途径。
\uses{lem:appO-lcopt, lem:appO-numerators, def:method-neighbourhoods}
\end{assumption}
\begin{proofsketch}测试。\end{proofsketch}

\begin{definition}[初态与末态中间态、取向与标准流]\label{def:appO-orientation}
（测试用合成陈述：内容从略。）
\uses{def:appO-lcopt, def:appO-leading-data}
\end{definition}

\section{第 2 节（测试）}\label{sec:appO-lcopt}

\begin{lemma}[光锥序中间态的取向与 $k^-$ 极点（精确形式） \tagR]\label{lem:appO-orientation}
（$R_{\mathrm{PT}}$：在重整化微扰论中对红外正规化的部分子态严格。）
（测试用合成陈述：内容从略。）
（本命题只涉及部分子态。）
\uses{def:appO-lcopt, lem:appO-lcopt, lem:appO-numerators, def:appO-orientation}
\end{lemma}

\begin{definition}[Glauber 区与 Glauber 夹挤]\label{def:appO-pinch}
（测试用合成陈述：内容从略。）
\uses{def:ir-glauber, lem:appO-orientation}
\end{definition}

\begin{lemma}[Glauber 幂次计数 \tagR]\label{lem:appO-glauber-pc}
（$R_{\mathrm{PT}}$：在重整化微扰论中对红外正规化的部分子态严格。）
（测试用合成陈述：内容从略。）
（本命题只涉及部分子态。）
% errtype: E-S
\uses{def:appO-leading-data, def:ir-glauber}
\end{lemma}

\begin{lemma}[硬顶点的光锥定域性与无权变量 $\mathcal E$ \tagR]\label{lem:appO-hard-locality}
（$R_{\mathrm{cond}}$：在所列编号假设之下证明完整；结论的地位不高于其中最弱的假设。）
前提：\cref{asm:appO-lcopt}。
（测试用合成陈述：内容从略。）
（本命题只涉及部分子态。）
% errtype: E-S
\uses{asm:appO-lcopt, lem:appO-lcopt, lem:appO-numerators, lem:appO-glauber-pc, def:appO-orientation}
\end{lemma}

\begin{definition}[末态链、割位置与约化权重]\label{def:appO-chain}
（测试用合成陈述：内容从略。）
\uses{def:appO-orientation, def:prelim-cut}
\end{definition}

\begin{lemma}[割位置恒等式（含插入）]\label{lem:appO-cut-position}
（测试用合成陈述：内容从略。）
\uses{def:appO-chain}
\end{lemma}

\begin{lemma}[光锥幺正性引理（闭时路径形式） \tagR]\label{lem:appO-unitarity}
（$R_{\mathrm{cond}}$：在所列编号假设之下证明完整；结论的地位不高于其中最弱的假设。）
前提：\cref{asm:appO-lcopt}。
（测试用合成陈述：内容从略。）
（本命题只涉及部分子态。）
% errtype: E-S
\uses{asm:appO-lcopt, lem:appO-hard-locality, lem:appO-cut-position, lem:appJ-cut-hermiticity}
\end{lemma}

\section{第 3 节（测试）}\label{sec:appO-orientation}

\begin{remark}[幺正性引理用到了什么、没有消去什么]\label{rem:appO-cut-sum-content}
（测试用合成陈述：内容从略。）
% uses: none
\end{remark}

\begin{definition}[可局域端点、尺度适配窗口与软角区]\label{def:appO-windows}
（测试用合成陈述：内容从略。）
\uses{def:method-neighbourhoods, def:appO-leading-data}
\end{definition}

\begin{lemma}[等光锥时刻插入]\label{lem:appO-equal-time}
（测试用合成陈述：内容从略。）
% uses: none
\end{lemma}

\begin{lemma}[光锥序中的插入：(ET1)(ET2) 与等时值 \tagR]\label{lem:appO-insertion-values}
（$R_{\mathrm{PT}}$：在重整化微扰论中对红外正规化的部分子态严格。）
（测试用合成陈述：内容从略。）
（本命题只涉及部分子态。）
% errtype: E-S
\uses{lem:appO-equal-time, lem:appO-orientation, lem:appO-glauber-pc, def:appO-windows, def:method-subtraction}
\end{lemma}

\begin{remark}[常数 $c$ 与正规化依赖性]\label{rem:appO-constant-c}
（测试用合成陈述：内容从略。）
% uses: none
\end{remark}

\begin{proposition}[一条旁观者–旁观者线的局域化与 Glauber 算符 \tagR]\label{prop:appO-ss-localize}
（$R_{\mathrm{cond}}$：在所列编号假设之下证明完整；结论的地位不高于其中最弱的假设。）
前提：\cref{asm:appO-lcopt}。
（测试用合成陈述：内容从略。）
（本命题只涉及部分子态。）
% errtype: E-F
\uses{asm:appO-lcopt, lem:appO-unitarity, lem:appO-equal-time, lem:appO-insertion-values, lem:appO-glauber-pc, def:appO-windows}
\end{proposition}

\begin{definition}[插入扇区、类 (SEC) 与结构条件 (EQ)(ADJ)(REAL)]\label{def:appO-sectors}
（测试用合成陈述：内容从略。）
\uses{def:appO-lcopt, def:appO-orientation}
\end{definition}

\section{第 4 节（测试）}\label{sec:appO-unitarity}

\begin{lemma}[非对易插入的联合局域化：扇区规则、次序权重与误差界 \tagR]\label{lem:appO-joint-localization}
（$R_{\mathrm{PT}}$：在重整化微扰论中对红外正规化的部分子态严格。）
（测试用合成陈述：内容从略。）
（本命题只涉及部分子态。）
% errtype: E-F
\uses{def:appO-sectors, lem:appO-equal-time, lem:appO-insertion-values}
\end{lemma}

\begin{lemma}[多极点外部因子的联合局域化 \tagR]\label{lem:appO-multipole-localization}
（$R_{\mathrm{PT}}$：在重整化微扰论中对红外正规化的部分子态严格。）
（测试用合成陈述：内容从略。）
（本命题只涉及部分子态。）
\uses{lem:appO-joint-localization}
\end{lemma}

\begin{lemma}[两条 SS 线的精确 $k^-$ 留数：象限公式 \tagR]\label{lem:appO-quadrant-residues}
（$R_{\mathrm{PT}}$：在重整化微扰论中对红外正规化的部分子态严格。）
（测试用合成陈述：内容从略。）
（本命题只涉及部分子态。）
\uses{lem:appO-joint-localization}
\end{lemma}

\begin{lemma}[远离壳的软线：在 $S$ 圈积分之前局域化 \tagR]\label{lem:appO-preintegration-localization}
（$R_{\mathrm{PT}}$：在重整化微扰论中对红外正规化的部分子态严格。）
（测试用合成陈述：内容从略。）
（本命题只涉及部分子态。）
\uses{lem:appO-equal-time, lem:appO-multipole-localization}
\end{lemma}

\begin{lemma}[Glauber 插入的非阿贝尔算符结构 \tagR]\label{lem:appO-glauber-operator}
（$R_{\mathrm{def}}$：由算符定义直接推出。）
（测试用合成陈述：内容从略。）
% uses: none
\end{lemma}

\begin{definition}[全局两时模型 $\mathbb M_{\rm glob}$、NF 域、模式与次序权重]\label{def:appO-glob-model}
（测试用合成陈述：内容从略。）
\uses{def:appO-sectors, lem:appO-glauber-pc, lem:appO-quadrant-residues, lem:appO-glauber-operator}
\end{definition}

\begin{definition}[正规子类 $\mathfrak X$ 与双侧窗口]\label{def:appO-regulator-class}
（测试用合成陈述：内容从略。）
\uses{def:appO-glob-model}
\end{definition}

\section{第 5 节（测试）}\label{sec:appO-localization}

\begin{lemma}[$k^+$ 围道的转动]\label{lem:appO-kplus-rotation}
（测试用合成陈述：内容从略。）
\uses{def:appO-glob-model, def:appO-regulator-class, lem:appO-quadrant-residues}
\end{lemma}

\begin{lemma}[射线恒等式]\label{lem:appO-ray-identity}
（测试用合成陈述：内容从略。）
% uses: none
\end{lemma}

\begin{lemma}[转动围道上双插入喷注函数的界]\label{lem:appO-two-insertion-bounds}
（测试用合成陈述：内容从略。）
\uses{def:appO-sectors, lem:appO-kplus-rotation}
\end{lemma}

\begin{lemma}[半重合对的色代数 \tagR]\label{lem:appO-halfcoinc-colour}
（$R_{\mathrm{def}}$：由算符定义直接推出。）
（测试用合成陈述：内容从略。）
\uses{lem:appO-glauber-operator}
\end{lemma}

\begin{definition}[旁观者交换族与积分映射]\label{def:appO-family}
（测试用合成陈述：内容从略。）
\uses{def:appO-leading-data, def:appO-windows}
\end{definition}

\begin{lemma}[与侧分配无关的窗口消去窗口边缘项 \tagR]\label{lem:appO-side-independent-windows}
（$R_{\mathrm{cond}}$：在所列编号假设之下证明完整；结论的地位不高于其中最弱的假设。）
前提：\cref{asm:appO-lcopt}。
（测试用合成陈述：内容从略。）
（本命题只涉及部分子态。）
\uses{asm:appO-lcopt, def:appO-family, def:appO-windows, prop:appO-ss-localize, lem:appO-orientation}
\end{lemma}

\begin{proposition}[旁观者–旁观者 Glauber 交换在族求和后相消 \tagR]\label{prop:appO-ss-cancel}
（$R_{\mathrm{cond}}$：在所列编号假设之下证明完整；结论的地位不高于其中最弱的假设。）
前提：\cref{asm:appO-lcopt}。
（测试用合成陈述：内容从略。）
（本命题只涉及部分子态。）
% errtype: E-F
\uses{asm:appO-lcopt, lem:appO-unitarity, prop:appO-ss-localize, lem:appO-joint-localization, lem:appO-glauber-operator, def:appO-family, lem:appO-side-independent-windows}
\end{proposition}

\section{第 6 节（测试）}\label{sec:appO-rewriting}

\begin{example}[单条 Glauber 交换]\label{ex:appO-single-exchange}
（测试用合成陈述：内容从略。）
\uses{lem:appO-equal-time, lem:appO-insertion-values, prop:appO-ss-localize}
\end{example}

\begin{example}[两条 Glauber 交换与 $p_{\rm opp}$]\label{ex:appO-two-exchange}
（测试用合成陈述：内容从略。）
\uses{lem:appO-glauber-operator, prop:dy-eikonal-twogluon}
\end{example}

\begin{assumption}[相邻交换对称性（ADJ） \tagP]\label{asm:appO-adjacent-swap}
（测试用合成陈述：内容从略。）
地位：测试状态句。
文献：测试。
被使用于：\cref{prop:appO-twoexch-spectator}、\cref{prop:appO-pergraph-weights}、\cref{prop:appO-halfcoinc-sensitivity}。
解除途径：测试解除途径。
\uses{def:appO-sectors}
\end{assumption}

\begin{assumption}[旁观者喷注函数的实性（REAL） \tagP]\label{asm:appO-real-jets}
（测试用合成陈述：内容从略。）
地位：测试状态句。
文献：测试。
被使用于：\cref{prop:appO-twoexch-spectator}、\cref{prop:appO-pergraph-weights}、\cref{prop:appO-halfcoinc-sensitivity}。
解除途径：测试解除途径。
\uses{def:appO-sectors}
\end{assumption}

\begin{assumption}[共线端点的匹配 \tagA]\label{asm:appO-collinear-ends}
（测试用合成陈述：内容从略。）
地位：测试状态句。
文献：测试。
被使用于：\cref{prop:appO-twoexch-spectator}、\cref{prop:appO-twoexch-regulated}、\cref{prop:appO-halfcoinc-sensitivity}。
解除途径：测试解除途径。
\uses{def:appO-regulator-class, def:appO-glob-model}
\end{assumption}
\begin{proofsketch}测试。\end{proofsketch}

\begin{proposition}[全重合 SS 对的 $O(g^4)$ 侧求和（未割线）：全局两时模型中的相消 \tagP]\label{prop:appO-twoexch-spectator}
（测试用合成陈述：内容从略。）
\uses{def:appO-glob-model, lem:appO-joint-localization, lem:appO-quadrant-residues, lem:appO-kplus-rotation, lem:appO-ray-identity, lem:appO-two-insertion-bounds, lem:appO-glauber-operator, asm:appO-lcopt, asm:appO-adjacent-swap, asm:appO-real-jets, asm:appO-collinear-ends}
\end{proposition}
\paragraph{地位与缺口}测试：本陈述的地位由\cref{def:appO-glob-model}、\cref{asm:appO-lcopt}、\cref{asm:appO-adjacent-swap}、\cref{asm:appO-real-jets}、\cref{prop:appO-strict-window-shape}、\cref{asm:appO-collinear-ends}、\cref{rem:appO-cut-ss}、\cref{asm:appO-mixing}、\cref{lem:appO-joint-localization}、\cref{lem:appO-quadrant-residues}、\cref{lem:appO-kplus-rotation}、\cref{lem:appO-ray-identity}、\cref{lem:appO-two-insertion-bounds}、\cref{lem:appO-glauber-operator}决定。

\begin{proposition}[正规化版本、端点极点与双侧正规子 \tagP]\label{prop:appO-twoexch-regulated}
（测试用合成陈述：内容从略。）
\uses{prop:appO-twoexch-spectator, def:appO-regulator-class, lem:appO-multipole-localization, lem:appO-ray-identity, asm:appO-collinear-ends}
\end{proposition}
\paragraph{地位与缺口}测试：本陈述的地位由\cref{prop:appO-twoexch-spectator}、\cref{asm:appO-collinear-ends}、\cref{prop:appO-strict-window-shape}、\cref{def:appO-regulator-class}、\cref{lem:appO-multipole-localization}、\cref{lem:appO-ray-identity}决定。

\section{第 7 节（测试）}\label{sec:appO-deformation}

\begin{proposition}[严格窗口及其对正规子形状的依赖 \tagR]\label{prop:appO-strict-window-shape}
（$R_{\mathrm{PT}}$：在重整化微扰论中对红外正规化的部分子态严格。）
（测试用合成陈述：内容从略。）
（本命题只涉及部分子态。）
\uses{lem:appO-joint-localization, lem:appO-multipole-localization, lem:appO-glauber-operator}
\end{proposition}

\begin{proposition}[全重合对的逐图次序权重不由窗口数据决定 \tagP]\label{prop:appO-pergraph-weights}
（测试用合成陈述：内容从略。）
\uses{def:appO-glob-model, lem:appO-quadrant-residues, prop:appO-twoexch-spectator, asm:appO-adjacent-swap, asm:appO-real-jets}
\end{proposition}
\paragraph{地位与缺口}测试：本陈述的地位由\cref{asm:appO-envelope}、\cref{def:appO-glob-model}、\cref{lem:appO-quadrant-residues}、\cref{prop:appO-twoexch-spectator}、\cref{asm:appO-adjacent-swap}、\cref{asm:appO-real-jets}决定。

\begin{remark}[半重合对：无正规子的两时模型]\label{rem:appO-halfcoinc}
（测试用合成陈述：内容从略。）
% uses: none
\end{remark}

\begin{proposition}[半重合对：$W_{12}$ 的实部与对 $k^-$ 正规子的敏感性 \tagP]\label{prop:appO-halfcoinc-sensitivity}
（测试用合成陈述：内容从略。）
\uses{def:appO-glob-model, def:appO-regulator-class, lem:appO-multipole-localization, lem:appO-kplus-rotation, lem:appO-two-insertion-bounds, lem:appO-halfcoinc-colour, asm:appO-adjacent-swap, asm:appO-real-jets, asm:appO-collinear-ends}
\end{proposition}
\paragraph{地位与缺口}测试：本陈述的地位由\cref{prop:appO-twoexch-spectator}、\cref{asm:appO-collinear-ends}、\cref{def:appO-glob-model}、\cref{def:appO-regulator-class}、\cref{lem:appO-multipole-localization}、\cref{lem:appO-kplus-rotation}、\cref{lem:appO-two-insertion-bounds}、\cref{lem:appO-halfcoinc-colour}、\cref{asm:appO-adjacent-swap}、\cref{asm:appO-real-jets}决定。

\begin{remark}[被割的直接 SS 线（实发射）]\label{rem:appO-cut-ss}
（测试用合成陈述：内容从略。）
% uses: none
\end{remark}

\begin{proposition}[$O(g^6)$：直接线梯形图类在幺正性下不封闭 \tagA]\label{prop:appO-g6-ladder}
（测试用合成陈述：内容从略。）
\uses{prop:dy-eikonal-twogluon, lem:appO-glauber-operator}
\end{proposition}
\begin{proofsketch}测试。\end{proofsketch}
\paragraph{地位与缺口}测试：本陈述的地位由\cref{prop:dy-eikonal-twogluon}、\cref{lem:appO-glauber-operator}、\cref{asm:appO-mixing}决定。

\begin{lemma}[整函数载子]\label{lem:appO-carrier}
（测试用合成陈述：内容从略。）
% uses: none
\end{lemma}

\section{第 8 节（测试）}\label{sec:appO-wilson}

\begin{definition}[测量相容的割图求和与弱等价]\label{def:appO-weak-equivalence}
（测试用合成陈述：内容从略。）
\uses{def:appO-leading-data, def:method-neighbourhoods, def:method-power-suppressed, def:method-subtraction}
\end{definition}

\begin{definition}[改写映射 $\mathfrak r_R$]\label{def:appO-rewriting}
（测试用合成陈述：内容从略。）
\uses{lem:appO-carrier, def:appO-windows, lem:appO-unitarity, lem:appO-insertion-values, lem:appO-joint-localization, def:method-subtraction}
\end{definition}

\begin{assumption}[纤维积分误差的可积包络 \tagA]\label{asm:appO-envelope}
（测试用合成陈述：内容从略。）
地位：测试状态句。
文献：测试。
被使用于：\cref{prop:method-glauber-deformation}、\cref{prop:dy-deformation}、\cref{thm:dy-factorization}、\cref{asm:appO-mixing}、\cref{prop:appO-one-sided}、\cref{thm:appO-main}、\cref{cor:appO-interface}、\cref{prop:appO-tmd}。
解除途径：测试解除途径。
\uses{def:appO-rewriting, def:appO-windows}
\end{assumption}
\begin{proofsketch}测试。\end{proofsketch}

\begin{assumption}[Glauber–软混合与重合连接的完整图集相消 \tagA]\label{asm:appO-mixing}
（测试用合成陈述：内容从略。）
地位：测试状态句。
文献：测试。
被使用于：\cref{prop:method-glauber-deformation}、\cref{prop:dy-deformation}、\cref{thm:dy-factorization}、\cref{lem:tmd-glauber}、\cref{thm:tmd-factorization}、\cref{prop:appO-one-sided}、\cref{thm:appO-main}、\cref{cor:appO-interface}、\cref{prop:appO-tmd}。
解除途径：测试解除途径。
\uses{def:appO-windows, def:appO-family, def:appO-weak-equivalence, asm:appO-envelope, lem:appO-unitarity, lem:appO-equal-time, lem:appO-preintegration-localization}
\end{assumption}
\begin{proofsketch}测试。\end{proofsketch}

\begin{proposition}[割图求和后的单侧性（弱等价） \tagR]\label{prop:appO-one-sided}
（$R_{\mathrm{cond}}$：在所列编号假设之下证明完整；结论的地位不高于其中最弱的假设。）
前提：\cref{asm:appO-lcopt}、\cref{asm:appO-envelope}、\cref{asm:appN-envelope}、\cref{asm:appN-separation}、\cref{asm:appO-mixing}。
（测试用合成陈述：内容从略。）
（本命题只涉及部分子态。）
% errtype: E-W
\uses{asm:appO-lcopt, asm:appO-envelope, asm:appN-envelope, asm:appN-separation, asm:appO-mixing, def:appO-rewriting, def:appO-weak-equivalence, lem:appO-unitarity, lem:appO-insertion-values, lem:appO-joint-localization, prop:appO-ss-localize, lem:appO-carrier, lem:appO-orientation}
\end{proposition}

\begin{remark}[为何只有弱等价；SA/AS 线；窗口外]\label{rem:appO-weak}
（测试用合成陈述：内容从略。）
% uses: none
\end{remark}

\begin{remark}[H-发射端点的障碍]\label{rem:appO-h-emission}
（测试用合成陈述：内容从略。）
% uses: none
\end{remark}

\section{第 9 节（测试）}\label{sec:appO-main}

\begin{definition}[形变映射与 Glauber 流]\label{def:appO-deformation-map}
（测试用合成陈述：内容从略。）
\uses{def:appO-windows, def:appO-leading-data, def:method-neighbourhoods}
\end{definition}

\begin{lemma}[形变不跨越奇点 \tagR]\label{lem:appO-no-crossing}
（$R_{\mathrm{PT}}$：在重整化微扰论中对红外正规化的部分子态严格。）
（测试用合成陈述：内容从略。）
（本命题只涉及部分子态。）
\uses{def:appO-deformation-map, def:appO-rewriting, lem:appO-carrier, lem:appO-orientation, def:appN-tmd-approx}
\end{lemma}

\begin{definition}[分级构型]\label{def:appO-graded}
（测试用合成陈述：内容从略。）
\uses{def:appO-deformation-map}
\end{definition}

\begin{definition}[形变目标与相容性]\label{def:appO-target}
（测试用合成陈述：内容从略。）
\uses{def:appO-windows, def:appO-weak-equivalence, def:appO-pinch, def:method-subtraction}
\end{definition}

\begin{proposition}[$\mathcal G_0$ 与分级构型的形变目标 \tagR]\label{prop:appO-target}
（$R_{\mathrm{PT}}$：在重整化微扰论中对红外正规化的部分子态严格。）
（测试用合成陈述：内容从略。）
（本命题只涉及部分子态。）
\uses{lem:appO-no-crossing, def:appO-graded, def:appO-target, def:appO-deformation-map}
\end{proposition}

\begin{assumption}[一般构型、相容性、软角区与 H-发射端点 \tagA]\label{asm:appO-deformation}
（测试用合成陈述：内容从略。）
地位：测试状态句。
文献：测试。
被使用于：\cref{prop:method-glauber-deformation}、\cref{prop:dy-deformation}、\cref{thm:dy-factorization}、\cref{lem:tmd-glauber}、\cref{prop:appN-dy}、\cref{thm:appO-main}、\cref{cor:appO-interface}、\cref{prop:appO-tmd}。
解除途径：测试解除途径。
\uses{def:appO-target, def:appO-rewriting, def:appO-graded, prop:appO-ss-cancel, lem:appO-side-independent-windows}
\end{assumption}
\begin{proofsketch}测试。\end{proofsketch}

\begin{proposition}[变形围道上的程函近似 \tagR]\label{prop:appO-eikonal}
（$R_{\mathrm{PT}}$：在重整化微扰论中对红外正规化的部分子态严格。）
（测试用合成陈述：内容从略。）
（本命题只涉及部分子态。）
% errtype: E-P
\uses{def:appO-deformation-map, prop:appO-target, lem:appO-no-crossing, lem:method-i0-compat}
\end{proposition}

\section{第 10 节（测试）}\label{sec:appO-checklist}

\begin{assumption}[改写形式与族求和形式的 Collins 条件 \tagA]\label{asm:appO-rewritten-cs}
（测试用合成陈述：内容从略。）
地位：测试状态句。
文献：测试。
被使用于：\cref{prop:method-glauber-deformation}、\cref{prop:dy-deformation}、\cref{thm:dy-factorization}、\cref{lem:tmd-glauber}、\cref{thm:tmd-factorization}、\cref{prop:appN-dy}、\cref{prop:appO-past-wilson}、\cref{thm:appO-main}、\cref{cor:appO-interface}、\cref{prop:appO-tmd}。
解除途径：测试解除途径。
\uses{def:method-collins-conditions, def:appO-rewriting, def:appO-family}
\end{assumption}
\begin{proofsketch}测试。\end{proofsketch}

\begin{proposition}[Ward 恒等式与指向过去的 Wilson 线 \tagA]\label{prop:appO-past-wilson}
（测试用合成陈述：内容从略。）
\uses{prop:appO-eikonal, prop:appO-ss-cancel, asm:appO-rewritten-cs, asm:appN-ward-product, lem:method-ward-multi, prop:method-soft-wilson, prop:method-collinear-wilson, prop:prelim-wilson-rules, lem:appO-no-crossing}
\end{proposition}
\begin{proofsketch}测试。\end{proofsketch}
\paragraph{地位与缺口}测试：本陈述的地位由\cref{prop:appO-eikonal}、\cref{prop:method-collinear-wilson}、\cref{prop:appO-ss-cancel}、\cref{asm:appO-rewritten-cs}、\cref{asm:appN-ward-product}、\cref{lem:method-ward-multi}、\cref{prop:method-soft-wilson}、\cref{prop:prelim-wilson-rules}、\cref{lem:appO-no-crossing}、\cref{ex:dy-future-vs-past}决定。

\begin{remark}[类时与类空的非光锥 Wilson 线]\label{rem:appO-timelike}
（测试用合成陈述：内容从略。）
% uses: none
\end{remark}

\begin{remark}[在壳线的两粒子幺正性（$O(g^4)$）]\label{rem:appO-onshell-unitarity}
（测试用合成陈述：内容从略。）
% uses: none
\end{remark}

\begin{theorem}[色单态单举 DY 中 Glauber 区相消的条件定理 \tagR]\label{thm:appO-main}
（$R_{\mathrm{cond}}$：在所列编号假设之下证明完整；结论的地位不高于其中最弱的假设。）
前提：\cref{asm:appO-lcopt}、\cref{asm:appO-envelope}、\cref{asm:appO-mixing}、\cref{asm:appO-deformation}、\cref{asm:appO-rewritten-cs}、\cref{asm:appN-envelope}、\cref{asm:appN-separation}、\cref{asm:appN-ward-product}、\cref{prop:appO-past-wilson}、\cref{prop:method-collinear-wilson}。
（测试用合成陈述：内容从略。）
（本命题只涉及部分子态。）
\uses{asm:appO-lcopt, asm:appO-envelope, asm:appO-mixing, asm:appO-deformation, asm:appO-rewritten-cs, asm:appN-envelope, asm:appN-separation, asm:appN-ward-product, prop:appO-past-wilson, lem:appO-hard-locality, lem:appO-orientation, lem:appO-unitarity, prop:appO-one-sided, prop:appO-ss-cancel, prop:appO-target, prop:appO-eikonal, lem:method-soft-unitarity, prop:method-collinear-wilson}
\end{theorem}

\begin{corollary}[与附录 N 的接口 \tagR]\label{cor:appO-interface}
（$R_{\mathrm{cond}}$：在所列编号假设之下证明完整；结论的地位不高于其中最弱的假设。）
前提：\cref{asm:appO-lcopt}、\cref{asm:appO-envelope}、\cref{asm:appO-mixing}、\cref{asm:appO-deformation}、\cref{asm:appO-rewritten-cs}、\cref{asm:appN-envelope}、\cref{asm:appN-separation}、\cref{asm:appN-ward-product}、\cref{prop:appO-past-wilson}。
（测试用合成陈述：内容从略。）
（本命题只涉及部分子态。）
\uses{thm:appO-main, asm:appO-lcopt, asm:appO-envelope, asm:appO-mixing, asm:appO-deformation, asm:appO-rewritten-cs, asm:appN-envelope, asm:appN-separation, asm:appN-ward-product, prop:appO-past-wilson, lem:method-i0-compat, lem:method-covering, def:appO-target}
\end{corollary}

\begin{proposition}[TMD 版本 \tagR]\label{prop:appO-tmd}
（$R_{\mathrm{cond}}$：在所列编号假设之下证明完整；结论的地位不高于其中最弱的假设。）
前提：\cref{asm:appO-lcopt}、\cref{asm:appO-envelope}、\cref{asm:appO-mixing}、\cref{asm:appO-deformation}、\cref{asm:appO-rewritten-cs}、\cref{asm:appN-envelope}、\cref{asm:appN-separation}、\cref{asm:appN-ward-product}、\cref{prop:appO-past-wilson}。
（测试用合成陈述：内容从略。）
（本命题只涉及部分子态。）
\uses{thm:appO-main, asm:appO-lcopt, asm:appO-envelope, asm:appO-mixing, asm:appO-deformation, asm:appO-rewritten-cs, asm:appN-envelope, asm:appN-separation, asm:appN-ward-product, prop:appO-past-wilson, def:appN-tmd-approx, def:tmd-unsub, def:tmd-soft, lem:appO-unitarity, prop:appO-ss-cancel}
\end{proposition}

\section{第 11 节（测试）}\label{sec:appO-neighbours}

\begin{remark}[测量相容性的统一判据]\label{rem:appO-measurement}
（测试用合成陈述：内容从略。）
% uses: none
\end{remark}

\begin{remark}[$pp\to h_1h_2X$ 在 (O-3) 失效]\label{rem:appO-hh}
（测试用合成陈述：内容从略。）
% uses: none
\end{remark}

\begin{remark}[快度间隙与衍射]\label{rem:appO-gaps}
（测试用合成陈述：内容从略。）
% uses: none
\end{remark}

\begin{remark}[Doria–Frenkel–Taylor]\label{rem:appO-dft}
（测试用合成陈述：内容从略。）
% uses: none
\end{remark}

\begin{remark}[双 Drell–Yan]\label{rem:appO-double-dy}
（测试用合成陈述：内容从略。）
% uses: none
\end{remark}

\begin{remark}[与 Glauber SCET 的对照]\label{rem:appO-rs-dictionary}
（测试用合成陈述：内容从略。）
% uses: none
\end{remark}

\begin{remark}[与其他全阶论证的比较；TMD 版本；有争议的提议]\label{rem:appO-literature-comparison}
（测试用合成陈述：内容从略。）
% uses: none
\end{remark}
